% Options for packages loaded elsewhere
\PassOptionsToPackage{unicode}{hyperref}
\PassOptionsToPackage{hyphens}{url}
\PassOptionsToPackage{dvipsnames,svgnames,x11names}{xcolor}
\documentclass[
  11pt,
]{article}
\usepackage{xcolor}
\usepackage[margin=2.5cm]{geometry}
\usepackage{amsmath,amssymb}
\setcounter{secnumdepth}{-\maxdimen} % remove section numbering
\usepackage{iftex}
\ifPDFTeX
  \usepackage[T1]{fontenc}
  \usepackage[utf8]{inputenc}
  \usepackage{textcomp} % provide euro and other symbols
\else % if luatex or xetex
  \usepackage{unicode-math} % this also loads fontspec
  \defaultfontfeatures{Scale=MatchLowercase}
  \defaultfontfeatures[\rmfamily]{Ligatures=TeX,Scale=1}
\fi
\usepackage{lmodern}
\ifPDFTeX\else
  % xetex/luatex font selection
    \setmainfont[]{Charter}
    \setmonofont[Scale=0.85]{Menlo}
  \setmathfont[]{STIX Two Math}
\fi
% Use upquote if available, for straight quotes in verbatim environments
\IfFileExists{upquote.sty}{\usepackage{upquote}}{}
\IfFileExists{microtype.sty}{% use microtype if available
  \usepackage[]{microtype}
  \UseMicrotypeSet[protrusion]{basicmath} % disable protrusion for tt fonts
}{}
\usepackage{setspace}
\makeatletter
\@ifundefined{KOMAClassName}{% if non-KOMA class
  \IfFileExists{parskip.sty}{%
    \usepackage{parskip}
  }{% else
    \setlength{\parindent}{0pt}
    \setlength{\parskip}{6pt plus 2pt minus 1pt}}
}{% if KOMA class
  \KOMAoptions{parskip=half}}
\makeatother
\usepackage{longtable,booktabs,array}
\usepackage{calc} % for calculating minipage widths
% Correct order of tables after \paragraph or \subparagraph
\usepackage{etoolbox}
\makeatletter
\patchcmd\longtable{\par}{\if@noskipsec\mbox{}\fi\par}{}{}
\makeatother
% Allow footnotes in longtable head/foot
\IfFileExists{footnotehyper.sty}{\usepackage{footnotehyper}}{\usepackage{footnote}}
\makesavenoteenv{longtable}
\ifLuaTeX
\usepackage[bidi=basic]{babel}
\else
\usepackage[bidi=default]{babel}
\fi
\babelprovide[main,import]{brazilian}
\ifPDFTeX
\else
\babelfont{rm}[]{Charter}
\fi
% get rid of language-specific shorthands (see #6817):
\let\LanguageShortHands\languageshorthands
\def\languageshorthands#1{}
\setlength{\emergencystretch}{3em} % prevent overfull lines
\providecommand{\tightlist}{%
  \setlength{\itemsep}{0pt}\setlength{\parskip}{0pt}}
\usepackage{bookmark}
\IfFileExists{xurl.sty}{\usepackage{xurl}}{} % add URL line breaks if available
\urlstyle{same}
\hypersetup{
  pdftitle={Guia de estudos},
  pdfauthor={Giovanni Vasconcelos · CESUPA Tech Summit 2026},
  pdflang={pt-BR},
  colorlinks=true,
  linkcolor={black},
  filecolor={Maroon},
  citecolor={Blue},
  urlcolor={blue},
  pdfcreator={LaTeX via pandoc}}

\title{Guia de estudos}
\usepackage{etoolbox}
\makeatletter
\providecommand{\subtitle}[1]{% add subtitle to \maketitle
  \apptocmd{\@title}{\par {\large #1 \par}}{}{}
}
\makeatother
\subtitle{Preparação para o workshop Attention is All You Need}
\author{Giovanni Vasconcelos · CESUPA Tech Summit 2026}
\date{}

\begin{document}
\maketitle

\renewcommand*\contentsname{Sumário}
{
\hypersetup{linkcolor=}
\setcounter{tocdepth}{3}
\tableofcontents
}
\setstretch{1.15}
\section{Como usar este guia}\label{como-usar-este-guia}

O guia supõe familiaridade com aprendizado de máquina e com a ideia
geral de um Transformer. O objetivo é outro: chegar ao workshop capaz de
derivar cada peça no quadro, justificar cada escolha de projeto e
responder perguntas fora do roteiro sem hesitar.

Cada módulo tem três partes: o que estudar (com fontes), os pontos que
você deve conseguir explicar e exercícios de verificação. Os exercícios
são a parte mais importante. Se você não consegue resolver um deles sem
consultar material, o módulo não está pronto.

\subsection{Cronograma sugerido (três
semanas)}\label{cronograma-sugerido-truxeas-semanas}

\begin{longtable}[]{@{}
  >{\raggedright\arraybackslash}p{(\linewidth - 4\tabcolsep) * \real{0.1053}}
  >{\raggedright\arraybackslash}p{(\linewidth - 4\tabcolsep) * \real{0.7018}}
  >{\raggedright\arraybackslash}p{(\linewidth - 4\tabcolsep) * \real{0.1930}}@{}}
\toprule\noalign{}
\begin{minipage}[b]{\linewidth}\raggedright
Semana
\end{minipage} & \begin{minipage}[b]{\linewidth}\raggedright
Módulos
\end{minipage} & \begin{minipage}[b]{\linewidth}\raggedright
Carga aproximada
\end{minipage} \\
\midrule\noalign{}
\endhead
\bottomrule\noalign{}
\endlastfoot
1 & 0 (revisão), 1 (artigo), 2 (atenção) & 8 a 10 h \\
2 & 3 (componentes), 4 (treino), 5 (geração), 8 (código do workshop) & 8
a 10 h \\
3 & 6 (do GPT aos LLMs), 7 (xadrez e modelos de mundo), 9 (perguntas
difíceis), ensaio & 6 a 8 h \\
\end{longtable}

Uma ordem de leitura que funciona bem: primeiro o vídeo de Karpathy
(Módulo 8), que dá a visão de conjunto em código; depois o artigo, lido
com o roteiro do Módulo 1; e só então os aprofundamentos.

\section{Módulo 0. Revisão de
pré-requisitos}\label{muxf3dulo-0.-revisuxe3o-de-pruxe9-requisitos}

\textbf{Estudar.}

\begin{itemize}
\tightlist
\item
  Álgebra linear: produto escalar como medida de similaridade,
  multiplicação de matrizes em termos de linhas e colunas, projeções
  lineares. Referência: Prince (2023), apêndice B; ou o capítulo 2 de
  Goodfellow, Bengio e Courville (2016).
\item
  Probabilidade: softmax, entropia cruzada, relação entre entropia
  cruzada e máxima verossimilhança. Goodfellow et al.~(2016), seções
  3.13 e 5.5.
\item
  Retropropagação e diferenciação automática. Goodfellow et al.~(2016),
  seção 6.5.
\item
  PyTorch: tensores, broadcasting, \texttt{nn.Module},
  \texttt{register\_buffer}, \texttt{torch.no\_grad}, o ciclo
  \texttt{zero\_grad}/\texttt{backward}/\texttt{step}.
\end{itemize}

\textbf{Você deve conseguir explicar.}

\begin{itemize}
\tightlist
\item
  Por que a entropia cruzada é a perda natural para classificação sobre
  um vocabulário.
\item
  A diferença entre \texttt{x.T} e \texttt{x.transpose(-2,\ -1)} em um
  tensor de três dimensões.
\item
  O que \texttt{keepdims=True} muda no broadcasting.
\end{itemize}

\textbf{Exercícios.}

\begin{enumerate}
\def\labelenumi{\arabic{enumi}.}
\tightlist
\item
  Mostre que \(\text{softmax}(z) = \text{softmax}(z - c)\) para qualquer
  escalar \(c\).
\item
  Mostre que o jacobiano da softmax é
  \(\partial p_i / \partial z_j = p_i(\delta_{ij} - p_j)\), isto é,
  \(J = \text{diag}(p) - pp^\top\). Conclua que, quando \(p\) se
  aproxima de um vetor one-hot, todas as entradas de \(J\) tendem a
  zero.
\item
  Um modelo que prevê a distribuição uniforme sobre 33 símbolos tem
  perda \(\ln 33 \approx 3{,}50\). Confira esse valor e relacione com a
  perda inicial impressa no treino de xadrez.
\end{enumerate}

\section{Módulo 1. Leitura guiada do
artigo}\label{muxf3dulo-1.-leitura-guiada-do-artigo}

\textbf{Estudar.} Vaswani et al.~(2017), \emph{Attention Is All You
Need}. Leia o artigo inteiro duas vezes: uma leitura corrida e uma
segunda com as perguntas abaixo. A nota
\texttt{2017-01-01\_vaswani-attention-is-all-you-need} do seu vault tem
um resumo útil para revisão.

\textbf{Roteiro por seção.}

\begin{itemize}
\tightlist
\item
  \textbf{Seções 1 e 2 (introdução e contexto).} Qual é o problema das
  RNNs que o artigo ataca? Qual era o papel da atenção antes de 2017
  (Bahdanau et al., 2015)? O que exatamente o título quer dizer com
  ``all you need''?
\item
  \textbf{Seção 3.1 (arquitetura).} Desenhe a Figura 1 de memória.
  Identifique as três formas de atenção: self-attention no encoder,
  self-attention mascarada no decoder e atenção encoder-decoder.
\item
  \textbf{Seção 3.2 (atenção).} Leia a nota de rodapé 4, que justifica o
  fator \(1/\sqrt{d_k}\). Compare atenção aditiva e multiplicativa: por
  que os autores preferem a multiplicativa?
\item
  \textbf{Seção 3.3 (FFN).} Por que \(d_{ff} = 4 \cdot d_{model}\)? (Não
  há justificativa teórica forte; é uma escolha empírica que se
  manteve.)
\item
  \textbf{Seção 3.4 (embeddings).} Os pesos da camada de embedding e da
  camada linear final são compartilhados, e os embeddings são
  multiplicados por \(\sqrt{d_{model}}\). Por quê?
\item
  \textbf{Seção 3.5 (codificação de posição).} Por que senos e cossenos?
  Qual o argumento sobre posições relativas?
\item
  \textbf{Seção 4 (por que self-attention).} Explique cada coluna da
  Tabela 1. Esta tabela é o slide 7 do workshop.
\item
  \textbf{Seção 5 (treino).} Anote os hiperparâmetros: Adam com
  \(\beta_1 = 0{,}9\), \(\beta_2 = 0{,}98\), \(\epsilon = 10^{-9}\);
  \emph{warmup} de 4000 passos; dropout 0,1; \emph{label smoothing} 0,1.
  O modelo base treinou 100 mil passos em 12 horas em 8 GPUs P100; o
  grande, 300 mil passos em 3,5 dias.
\item
  \textbf{Seção 6 (resultados).} 28,4 BLEU em inglês-alemão e 41,8 em
  inglês-francês (WMT 2014). Leia a Tabela 3 com atenção: ela é um
  estudo de ablação e responde a várias perguntas prováveis da plateia
  (efeito do número de cabeças, de \(d_k\), de dropout, de embeddings de
  posição aprendidos).
\end{itemize}

\textbf{Você deve conseguir explicar.}

\begin{itemize}
\tightlist
\item
  O fluxo de dados completo do modelo de tradução, da frase em inglês à
  distribuição sobre a próxima palavra em alemão.
\item
  A diferença entre as três formas de atenção e de onde vêm Q, K e V em
  cada uma.
\item
  A taxa de aprendizado do artigo:
  \(\text{lr} = d_{model}^{-0,5} \cdot \min(\text{step}^{-0,5},\ \text{step} \cdot \text{warmup}^{-1,5})\).
  Esboce a curva.
\end{itemize}

\textbf{Exercícios.}

\begin{enumerate}
\def\labelenumi{\arabic{enumi}.}
\tightlist
\item
  Conte os parâmetros do modelo base (\(N = 6\), \(d_{model} = 512\),
  \(d_{ff} = 2048\), vocabulário compartilhado de cerca de 37 mil
  tokens) e confirme que o total fica próximo de 65 milhões.
\item
  Segundo a Tabela 3, o que acontece com o BLEU quando se usa uma única
  cabeça? E com 32 cabeças?
\end{enumerate}

\section{Módulo 2. Atenção em
profundidade}\label{muxf3dulo-2.-atenuxe7uxe3o-em-profundidade}

\textbf{Estudar.}

\begin{itemize}
\tightlist
\item
  Prince (2023), capítulo 12, seções 12.1 a 12.4.
\item
  Alammar (2018), \emph{The Illustrated Transformer}.
\item
  3Blue1Brown (2024), capítulos 5 e 6 da série sobre redes neurais
  (\emph{Transformers} e \emph{Attention in transformers, visually
  explained}). Excelente fonte de analogias visuais para a exposição.
\item
  Bahdanau, Cho e Bengio (2015), para o contexto histórico.
\end{itemize}

\textbf{Pontos centrais.}

\emph{Atenção como média ponderada.} A saída na posição \(i\) é
\(\sum_j \alpha_{ij} v_j\), com
\(\alpha_{ij} = \text{softmax}_j(q_i \cdot k_j / \sqrt{d_k})\). É um
estimador de Nadaraya--Watson com núcleo exponencial: uma regressão por
núcleo em que a similaridade é aprendida. Essa leitura ajuda a explicar
a analogia do dicionário do slide 8.

\emph{O fator de escala.} Se \(q_i, k_i\) são independentes, com média 0
e variância 1, então \(\text{Var}(q_i k_i) = 1\) e
\(\text{Var}(q \cdot k) = d_k\). O desvio-padrão dos scores cresce como
\(\sqrt{d_k}\); a softmax satura e, pelo Exercício 2 do Módulo 0, os
gradientes desaparecem. Dividir por \(\sqrt{d_k}\) restaura a variância
unitária. O experimento da seção 1.3 do notebook mede isso: com 10 keys
e \(d_k = 512\), o maior peso médio sem escala é 0,95, e com escala,
0,31.

\emph{A máscara causal.} Somar \(-\infty\) antes da softmax zera
exatamente os pesos proibidos e renormaliza os restantes. A consequência
importante é o paralelismo no treino: uma passada pela sequência produz
\(T\) previsões ao mesmo tempo, todas com o contexto correto. Isso é
chamado de \emph{teacher forcing}, e é a razão principal de o treino de
Transformers ser tão mais eficiente que o de RNNs.

\emph{Por que várias cabeças.} Uma única softmax por posição produz uma
única média ponderada. Com \(h\) cabeças, a mesma posição pode agregar
informação de lugares diferentes, por critérios diferentes. Como cada
cabeça tem dimensão \(d_{model}/h\), o custo total é praticamente o
mesmo de uma cabeça de dimensão cheia. Existe um argumento de posto:
\(QK^\top\) tem posto no máximo \(d_k\), então cada cabeça é um mapa de
atenção de posto baixo.

\emph{Complexidade.} Tempo \(O(T^2 d)\) e memória \(O(T^2)\) por cabeça,
por causa da matriz de scores. É isso que limita o tamanho do contexto.
FlashAttention (Dao et al., 2022) calcula a atenção exata sem
materializar a matriz \(T \times T\) na memória da GPU, reorganizando o
cálculo em blocos.

\emph{Invariância a permutações.} Sem codificação de posição, permutar a
entrada apenas permuta a saída (a atenção é equivariante a permutações).
Sem a informação de ordem, ``o rei protege a rainha'' e ``a rainha
protege o rei'' teriam representações idênticas, a menos da ordem das
linhas.

\textbf{Exercícios.}

\begin{enumerate}
\def\labelenumi{\arabic{enumi}.}
\tightlist
\item
  Implemente a atenção em NumPy sem consultar o notebook, incluindo a
  máscara. Compare com \texttt{checagens.\_atencao\_ref}.
\item
  Demonstre formalmente a equivariância a permutações: para uma matriz
  de permutação \(P\),
  \(\text{Attention}(PX) = P \cdot \text{Attention}(X)\) quando não há
  máscara nem codificação de posição.
\item
  O que acontece com os pesos de atenção se todas as keys forem iguais?
  E se a query for o vetor nulo?
\item
  Se dividirmos por \(d_k\) em vez de \(\sqrt{d_k}\), qual o efeito
  sobre a distribuição dos pesos para \(d_k\) grande?
\end{enumerate}

\section{Módulo 3. Os demais
componentes}\label{muxf3dulo-3.-os-demais-componentes}

\textbf{Estudar.}

\begin{itemize}
\tightlist
\item
  He et al.~(2016), para conexões residuais; Ba, Kiros e Hinton (2016),
  para LayerNorm.
\item
  Xiong et al.~(2020), \emph{On Layer Normalization in the Transformer
  Architecture}: a comparação entre pre-LN e post-LN.
\item
  Elhage et al.~(2021), \emph{A Mathematical Framework for Transformer
  Circuits}: a leitura do Transformer como um \emph{residual stream} em
  que cada camada lê e escreve.
\item
  Geva et al.~(2021), \emph{Transformer Feed-Forward Layers Are
  Key-Value Memories}.
\item
  Su et al.~(2021), RoPE; Press, Smith e Lewis (2022), ALiBi.
\end{itemize}

\textbf{Pontos centrais.}

\emph{Conexões residuais.} Cada subcamada aprende uma correção
\(x \mapsto x + f(x)\), e o gradiente tem um caminho direto até as
camadas iniciais. Na visão de Elhage et al., o vetor de cada posição é
um canal de comunicação compartilhado: a atenção move informação entre
posições e a FFN processa informação dentro de uma posição.

\emph{Pre-LN e post-LN.} O artigo usa post-LN,
\(x = \text{LN}(x + f(x))\). O workshop usa pre-LN,
\(x = x + f(\text{LN}(x))\), como o GPT-2 e quase todos os modelos
atuais. Xiong et al.~(2020) mostram que o post-LN produz gradientes
grandes perto da saída no início do treino, o que exige \emph{warmup}
cuidadoso; o pre-LN treina de forma estável sem isso. Esta é uma
pergunta provável de quem tiver lido o artigo.

\emph{FFN.} Duas camadas lineares com ativação entre elas, aplicadas a
cada posição de forma independente. Geva et al.~(2021) interpretam a
primeira camada como um conjunto de chaves que detectam padrões e a
segunda como os valores correspondentes, o que aproxima a FFN de uma
memória associativa. A maior parte dos parâmetros de um Transformer está
nas FFNs.

\emph{Codificação de posição.} Três famílias. Absoluta senoidal (o
artigo); absoluta aprendida (GPT-2 e o workshop); relativa, em que a
posição entra no cálculo da atenção (RoPE, que rotaciona queries e keys
conforme a posição e é usada no Llama; ALiBi, que soma uma penalidade
proporcional à distância). As relativas generalizam melhor para
contextos maiores que os vistos no treino.

\emph{Contagem de parâmetros.} Por bloco, ignorando vieses e LayerNorm:
\(4d^2\) na atenção (\(W_Q, W_K, W_V, W_O\)) e \(8d^2\) na FFN (duas
matrizes \(d \times 4d\)), total \(12d^2\). O modelo da turma
(\(d = 128\), 4 blocos) tem \(12 \cdot 128^2 \cdot 4 \approx 786\) mil
parâmetros nos blocos, mais cerca de 40 mil nos embeddings: 0,83 milhão,
como o notebook imprime. O de referência (\(d = 256\), 6 blocos) tem
cerca de 4,8 milhões.

\textbf{Exercícios.}

\begin{enumerate}
\def\labelenumi{\arabic{enumi}.}
\tightlist
\item
  Refaça a contagem de parâmetros do modelo da turma incluindo vieses e
  LayerNorms e compare com \texttt{modelo.n\_params()}.
\item
  Implemente a codificação senoidal como alternativa ao
  \texttt{pos\_emb} (é o primeiro exercício para casa do notebook) e
  compare as curvas de perda.
\item
  Explique por que, com codificação senoidal, \(PE_{pos+k}\) é uma
  função linear de \(PE_{pos}\).
\end{enumerate}

\section{Módulo 4. Treino}\label{muxf3dulo-4.-treino}

\textbf{Estudar.}

\begin{itemize}
\tightlist
\item
  Loshchilov e Hutter (2019), AdamW.
\item
  Szegedy et al.~(2016), \emph{label smoothing}.
\item
  Micikevicius et al.~(2018), treino com precisão mista.
\end{itemize}

\textbf{Pontos centrais.}

\begin{itemize}
\tightlist
\item
  O objetivo é a entropia cruzada média do próximo token,
  \(\mathcal{L} = -\frac{1}{T}\sum_t \log p_\theta(x_{t+1} \mid x_{\le t})\).
  Minimizá-la equivale a maximizar a verossimilhança dos dados.
\item
  AdamW separa o \emph{weight decay} da atualização adaptativa do Adam,
  o que regulariza de forma mais previsível.
\item
  \emph{Warmup} seguido de decaimento (no workshop, cosseno no
  \texttt{pretreinar.py}; no artigo, inverso da raiz quadrada). O
  \emph{warmup} evita atualizações grandes enquanto as estimativas de
  segundo momento do Adam ainda são ruins.
\item
  Precisão mista: as multiplicações de matrizes rodam em fp16 e os pesos
  mestres ficam em fp32. O \texttt{GradScaler} multiplica a perda por um
  fator grande antes do \emph{backward}, para que gradientes pequenos
  não virem zero em fp16.
\item
  \emph{Gradient clipping} (usado no pré-treino) limita a norma do
  gradiente e evita passos catastróficos.
\end{itemize}

\textbf{Exercícios.}

\begin{enumerate}
\def\labelenumi{\arabic{enumi}.}
\tightlist
\item
  No notebook, a tarefa de ordenação usa alvo \(-1\) nas posições da
  entrada. Explique por que, e o que aconteceria se essas posições
  entrassem na perda.
\item
  Por que a perda de validação da Parte 3 fica próxima da perda de
  treino? O que você esperaria ver se o modelo fosse muito maior e
  treinado por muito mais tempo?
\end{enumerate}

\section{Módulo 5. Geração}\label{muxf3dulo-5.-gerauxe7uxe3o}

\textbf{Estudar.}

\begin{itemize}
\tightlist
\item
  Holtzman et al.~(2020), \emph{The Curious Case of Neural Text
  Degeneration} (amostragem por núcleo, ou \emph{top-p}).
\item
  Pope et al.~(2023), sobre inferência eficiente e KV cache (leitura
  opcional; o conceito basta).
\end{itemize}

\textbf{Pontos centrais.}

\begin{itemize}
\tightlist
\item
  Geração autorregressiva: amostra-se um token, ele é anexado ao
  contexto, e o modelo é executado de novo.
\item
  Temperatura \(\tau\): divide os logits antes da softmax.
  \(\tau \to 0\) aproxima a escolha gulosa; \(\tau\) grande aproxima a
  distribuição uniforme.
\item
  \emph{Top-k} e \emph{top-p}: truncam a distribuição antes de amostrar,
  eliminando a cauda de tokens improváveis. No xadrez, isso reduziria
  lances ilegais.
\item
  KV cache: durante a geração, as keys e values das posições anteriores
  não mudam, então podem ser guardadas. Cada novo token custa
  \(O(T \cdot d)\) em vez de recalcular tudo. A implementação do
  workshop não usa cache (por simplicidade), e por isso recalcula a
  sequência inteira a cada token.
\end{itemize}

\textbf{Exercícios.}

\begin{enumerate}
\def\labelenumi{\arabic{enumi}.}
\tightlist
\item
  Implemente \emph{top-k} no método \texttt{generate} e compare a taxa
  de legalidade do modelo da turma com \(k = 3\).
\item
  Calcule quantas operações a geração de 100 tokens custa com e sem KV
  cache, em função de \(T\).
\end{enumerate}

\section{Módulo 6. Do GPT aos grandes modelos de
linguagem}\label{muxf3dulo-6.-do-gpt-aos-grandes-modelos-de-linguagem}

\textbf{Estudar.}

\begin{itemize}
\tightlist
\item
  Radford et al.~(2018, 2019), GPT e GPT-2; Brown et al.~(2020), GPT-3 e
  o aprendizado em contexto.
\item
  Devlin et al.~(2019), BERT, para o contraste com os modelos
  encoder-only.
\item
  Sennrich, Haddow e Birch (2016), tokenização por BPE.
\item
  Kaplan et al.~(2020) e Hoffmann et al.~(2022), leis de escala. A
  conclusão de Hoffmann et al.~(o modelo ``Chinchilla'') é que, para um
  orçamento fixo de computação, o número de tokens de treino deve
  crescer na mesma proporção que o número de parâmetros, em torno de 20
  tokens por parâmetro.
\item
  Ouyang et al.~(2022), InstructGPT: ajuste por instruções e aprendizado
  por reforço com feedback humano (RLHF).
\end{itemize}

\textbf{Pontos centrais.}

\begin{itemize}
\tightlist
\item
  O modelo decoder-only venceu porque o objetivo ``prever o próximo
  token'' transforma qualquer texto em dado de treino, sem rotulação, e
  escala de forma previsível.
\item
  Um modelo pré-treinado é um completador de texto. O comportamento de
  assistente vem de etapas posteriores: ajuste supervisionado com
  exemplos de instruções e respostas, seguido de otimização por
  preferências humanas ou por recompensas verificáveis.
\item
  Diferenças arquiteturais comuns hoje em relação ao workshop: RMSNorm
  no lugar de LayerNorm, ativação SwiGLU na FFN, RoPE,
  \emph{grouped-query attention} (várias cabeças de query compartilham
  keys e values, o que reduz o KV cache) e, em alguns modelos,
  \emph{mixture of experts} na FFN.
\end{itemize}

\textbf{Exercício.} Monte uma tabela com três linhas (GPT-2, GPT-3, um
modelo aberto recente da família Llama) e colunas para número de
parâmetros, número de camadas, \(d_{model}\), número de cabeças, tamanho
do contexto e tokens de treino. Use os artigos ou os \emph{model cards}
originais.

\section{Módulo 7. Xadrez e modelos de
mundo}\label{muxf3dulo-7.-xadrez-e-modelos-de-mundo}

\textbf{Estudar.}

\begin{itemize}
\tightlist
\item
  Karvonen (2024), \emph{Emergent World Models and Latent Variable
  Estimation in Chess-Playing Language Models}, arXiv:2403.15498. É a
  referência direta da Parte 3: o formato de dados do workshop é o dele.
\item
  Li et al.~(2023), \emph{Emergent World Representations}: o experimento
  Othello-GPT.
\item
  Nanda, Lee e Wattenberg (2023), \emph{Emergent Linear Representations
  in World Models of Self-Supervised Sequence Models}, que mostra que a
  representação do tabuleiro no Othello-GPT é linear quando expressa
  como ``minhas peças'' e ``peças do adversário''.
\item
  Alain e Bengio (2016), sobre sondas lineares (\emph{linear probes})
  como método.
\end{itemize}

\textbf{Pontos centrais.}

\begin{itemize}
\tightlist
\item
  Uma sonda linear é um classificador linear treinado sobre as ativações
  congeladas do modelo para prever uma propriedade (por exemplo, ``que
  peça está em e4''). Se uma sonda linear acerta, a informação está
  presente e é linearmente acessível naquela camada.
\item
  Karvonen mostra que modelos treinados apenas com PGN desenvolvem
  representações do tabuleiro e até de uma variável latente de
  habilidade do jogador, e que intervir nessas representações altera o
  comportamento do modelo.
\item
  A taxa de legalidade do workshop é uma métrica comportamental, mais
  fraca que uma sonda: ela mostra que o modelo age como se soubesse a
  posição, sem mostrar onde nem como essa informação está representada.
\item
  Limitações a reconhecer: o modelo da turma é pequeno e treinado por
  minutos. Ele acerta aberturas porque elas são frequentes nos dados, e
  degrada no meio-jogo, onde as posições são novas. Parte do que parece
  compreensão é memorização de sequências comuns.
\end{itemize}

\textbf{Exercícios.}

\begin{enumerate}
\def\labelenumi{\arabic{enumi}.}
\tightlist
\item
  Esboce como você treinaria uma sonda linear para a casa e4 usando o
  modelo de referência: que ativações coletar, qual o alvo, como separar
  treino e teste.
\item
  Explique por que o formato de caracteres dificulta a sonda (um lance
  ocupa vários tokens) e como Karvonen lida com isso (as sondas são
  aplicadas na posição do ponto que precede cada lance das brancas).
\end{enumerate}

\section{Módulo 8. Domínio do código do
workshop}\label{muxf3dulo-8.-domuxednio-do-cuxf3digo-do-workshop}

\textbf{Estudar.}

\begin{itemize}
\tightlist
\item
  Karpathy (2023), \emph{Let's build GPT: from scratch, in code, spelled
  out} (vídeo, cerca de 2 horas) e o repositório \texttt{nanoGPT}. O
  código do workshop segue o mesmo estilo, com nomes em português.
\item
  Os arquivos do repositório: \texttt{mini\_gpt.py},
  \texttt{checagens.py}, \texttt{scripts/gerar\_notebooks.py},
  \texttt{scripts/pretreinar.py}, \texttt{scripts/preparar\_dados.py}.
\end{itemize}

\textbf{Você deve conseguir, sem consulta.}

\begin{itemize}
\tightlist
\item
  Escrever \texttt{Head}, \texttt{MultiHeadAttention} e \texttt{Block}
  do zero.
\item
  Explicar cada linha de \texttt{GPT.forward} com o shape de cada
  tensor.
\item
  Explicar como \texttt{DadosXadrez.batch} monta os exemplos (toda
  janela começa em um \texttt{;}, isto é, no início de uma partida) e
  por que isso importa para o modelo aprender a notação dos números de
  lance.
\item
  Explicar como \texttt{lance\_do\_modelo} monta o prompt (histórico em
  PGN terminado em \texttt{N.} ou em espaço, conforme a vez) e o que
  acontece quando o modelo erra.
\item
  Explicar a diferença entre o modelo da turma (4 blocos, 4 cabeças,
  \(d = 128\), contexto 256) e o de referência (6 blocos, 8 cabeças,
  \(d = 256\), contexto 384).
\end{itemize}

\textbf{Exercícios.}

\begin{enumerate}
\def\labelenumi{\arabic{enumi}.}
\tightlist
\item
  Resolva o \texttt{workshop\_aluno.ipynb} inteiro sem consultar o
  gabarito e cronometre.
\item
  Introduza de propósito cada erro da tabela da seção 5 do plano e
  observe a mensagem das verificações. Assim você reconhece o sintoma na
  sala.
\item
  Execute o notebook de pré-treino e anote a taxa de legalidade do
  modelo de referência.
\end{enumerate}

\section{Módulo 9. Perguntas difíceis da
plateia}\label{muxf3dulo-9.-perguntas-difuxedceis-da-plateia}

Respostas curtas, para ter prontas.

\textbf{``Por que se chama atenção?''} O nome vem da tradução automática
neural (Bahdanau et al., 2015): ao gerar cada palavra, o modelo ``presta
atenção'' a partes diferentes da frase de origem. Matematicamente, é uma
média ponderada com pesos que dependem dos dados.

\textbf{``Por que o custo é quadrático e isso importa?''} Cada posição
se compara com todas as outras: \(T^2\) comparações por cabeça e por
camada. Para contextos de centenas de milhares de tokens, isso domina o
custo; daí FlashAttention, atenção esparsa e alternativas como modelos
de espaço de estados.

\textbf{``O modelo entende xadrez?''} Depende do que se entende por
entender. Há evidência de que modelos desse tipo representam
internamente o estado do tabuleiro (Karvonen, 2024). Por outro lado, o
modelo da turma não planeja, não calcula variantes e erra regras no
meio-jogo. É uma boa oportunidade para distinguir comportamento,
representação e compreensão.

\textbf{``Qual a diferença para o AlphaZero ou o Stockfish?''} O
AlphaZero recebe o tabuleiro como entrada, combina uma rede com busca em
árvore (MCTS) e aprende por reforço jogando contra si mesmo. O Stockfish
usa busca alfa-beta com uma função de avaliação. O modelo do workshop
não faz busca nenhuma e nunca vê o tabuleiro: apenas imita partidas
humanas.

\textbf{``Por que caracteres e não palavras ou lances inteiros?''}
Simplicidade: 33 símbolos e nenhum tokenizador para explicar. O custo é
que cada lance vira vários tokens e as sequências ficam longas. LLMs
usam BPE, que é um meio-termo entre caracteres e palavras.

\textbf{``Os pesos de atenção explicam o que o modelo faz?''} Só
parcialmente. Jain e Wallace (2019) mostram que pesos de atenção
diferentes podem produzir a mesma saída. Os pesos dizem de onde a
informação é lida, não o que é feito com ela.

\textbf{``Por que dividir por \(\sqrt{d_k}\) e não normalizar de outro
jeito?''} É a correção mínima que torna a variância dos scores
independente da dimensão, supondo componentes independentes com
variância unitária. Variantes modernas às vezes normalizam queries e
keys com LayerNorm (\emph{QK-norm}), com o mesmo objetivo.

\textbf{``O que é o KV cache?''} Na geração, as keys e values das
posições passadas não mudam, então são armazenadas e reutilizadas. É o
principal consumidor de memória na inferência de LLMs.

\textbf{``Como o GPT passou a seguir instruções?''} Ajuste
supervisionado com pares de instrução e resposta, seguido de aprendizado
por reforço com preferências humanas (Ouyang et al., 2022). A
arquitetura não muda.

\textbf{``O que é o contexto de um modelo?''} O número máximo de tokens
que ele processa de uma vez (\texttt{block\_size} no workshop). No
código, \texttt{generate} corta o histórico para os últimos
\texttt{block\_size} tokens; em partidas muito longas, o modelo perde o
início da partida.

\textbf{``Transformers só servem para texto?''} Não. Vision Transformers
dividem imagens em blocos e tratam cada bloco como um token (Dosovitskiy
et al., 2021). O mesmo vale para áudio, proteínas e, como no workshop,
partidas de xadrez.

\textbf{``Por que o encoder não usa máscara?''} Porque na tradução a
frase de origem inteira está disponível desde o início. A máscara só é
necessária quando o modelo gera a sequência e não pode ver o futuro.

\section{Lista de autoverificação}\label{lista-de-autoverificauxe7uxe3o}

Antes do ensaio (D − 7), você deve conseguir, sem consulta:

\begin{itemize}
\tightlist
\item
  Derivar a variância de \(q \cdot k\) e explicar a saturação da softmax
  com o jacobiano.
\item
  Desenhar a Figura 1 do artigo e explicar as três formas de atenção.
\item
  Explicar a Tabela 1 do artigo coluna por coluna.
\item
  Escrever \texttt{Head}, \texttt{MultiHeadAttention} e \texttt{Block}
  em PyTorch sem erros de shape.
\item
  Contar os parâmetros do modelo da turma.
\item
  Explicar pre-LN e post-LN e por que o workshop diverge do artigo.
\item
  Explicar por que a máscara causal permite treinar \(T\) previsões em
  paralelo.
\item
  Explicar temperatura, \emph{top-k}, \emph{top-p} e KV cache.
\item
  Explicar o que é uma sonda linear e o resultado de Karvonen (2024).
\item
  Responder às perguntas do Módulo 9 em menos de um minuto cada.
\end{itemize}

\section{Referências}\label{referuxeancias}

\begin{itemize}
\tightlist
\item
  Alain, G.; Bengio, Y. (2016). Understanding intermediate layers using
  linear classifier probes. arXiv:1610.01644.
\item
  Alammar, J. (2018). The Illustrated Transformer.
  jalammar.github.io/illustrated-transformer.
\item
  Ba, J. L.; Kiros, J. R.; Hinton, G. E. (2016). Layer Normalization.
  arXiv:1607.06450.
\item
  Bahdanau, D.; Cho, K.; Bengio, Y. (2015). Neural Machine Translation
  by Jointly Learning to Align and Translate. ICLR.
\item
  Brown, T. et al.~(2020). Language Models are Few-Shot Learners.
  NeurIPS.
\item
  Dao, T. et al.~(2022). FlashAttention: Fast and Memory-Efficient Exact
  Attention with IO-Awareness. NeurIPS.
\item
  Devlin, J. et al.~(2019). BERT: Pre-training of Deep Bidirectional
  Transformers for Language Understanding. NAACL.
\item
  Dosovitskiy, A. et al.~(2021). An Image is Worth 16x16 Words:
  Transformers for Image Recognition at Scale. ICLR.
\item
  Elhage, N. et al.~(2021). A Mathematical Framework for Transformer
  Circuits. Transformer Circuits Thread.
\item
  Geva, M. et al.~(2021). Transformer Feed-Forward Layers Are Key-Value
  Memories. EMNLP.
\item
  Goodfellow, I.; Bengio, Y.; Courville, A. (2016). Deep Learning. MIT
  Press.
\item
  He, K. et al.~(2016). Deep Residual Learning for Image Recognition.
  CVPR.
\item
  Hoffmann, J. et al.~(2022). Training Compute-Optimal Large Language
  Models. NeurIPS.
\item
  Holtzman, A. et al.~(2020). The Curious Case of Neural Text
  Degeneration. ICLR.
\item
  Jain, S.; Wallace, B. C. (2019). Attention is not Explanation. NAACL.
\item
  Kaplan, J. et al.~(2020). Scaling Laws for Neural Language Models.
  arXiv:2001.08361.
\item
  Karpathy, A. (2023). Let's build GPT: from scratch, in code, spelled
  out. Vídeo (YouTube) e repositório nanoGPT.
\item
  Karvonen, A. (2024). Emergent World Models and Latent Variable
  Estimation in Chess-Playing Language Models. arXiv:2403.15498.
\item
  Li, K. et al.~(2023). Emergent World Representations: Exploring a
  Sequence Model Trained on a Synthetic Task. ICLR.
\item
  Loshchilov, I.; Hutter, F. (2019). Decoupled Weight Decay
  Regularization. ICLR.
\item
  Micikevicius, P. et al.~(2018). Mixed Precision Training. ICLR.
\item
  Nanda, N.; Lee, A.; Wattenberg, M. (2023). Emergent Linear
  Representations in World Models of Self-Supervised Sequence Models.
  BlackboxNLP.
\item
  Ouyang, L. et al.~(2022). Training language models to follow
  instructions with human feedback. NeurIPS.
\item
  Pope, R. et al.~(2023). Efficiently Scaling Transformer Inference.
  MLSys.
\item
  Press, O.; Smith, N. A.; Lewis, M. (2022). Train Short, Test Long:
  Attention with Linear Biases Enables Input Length Extrapolation. ICLR.
\item
  Prince, S. J. D. (2023). Understanding Deep Learning. MIT Press.
\item
  Radford, A. et al.~(2018). Improving Language Understanding by
  Generative Pre-Training. OpenAI.
\item
  Radford, A. et al.~(2019). Language Models are Unsupervised Multitask
  Learners. OpenAI.
\item
  Sennrich, R.; Haddow, B.; Birch, A. (2016). Neural Machine Translation
  of Rare Words with Subword Units. ACL.
\item
  Su, J. et al.~(2021). RoFormer: Enhanced Transformer with Rotary
  Position Embedding. arXiv:2104.09864.
\item
  Szegedy, C. et al.~(2016). Rethinking the Inception Architecture for
  Computer Vision. CVPR.
\item
  Vaswani, A. et al.~(2017). Attention Is All You Need. NeurIPS.
  arXiv:1706.03762.
\item
  Xiong, R. et al.~(2020). On Layer Normalization in the Transformer
  Architecture. ICML.
\item
  3Blue1Brown (2024). Neural Networks, capítulos 5 e 6. Vídeos
  (YouTube).
\end{itemize}

\end{document}
